\documentclass[11pt]{article}
\usepackage[margin=1in]{geometry}
\usepackage[T1]{fontenc}
\usepackage[utf8]{inputenc}
\usepackage{booktabs}
\usepackage{array}
\usepackage{hyperref}
\usepackage{enumitem}

\hypersetup{
  colorlinks=true,
  urlcolor=blue,
  linkcolor=blue,
  citecolor=blue
}

\setlength{\parindent}{0pt}
\setlength{\parskip}{0.6em}
\renewcommand{\arraystretch}{1.15}

\title{When Surgical Edits Leak:\\Midterm Progress Report}
\author{Matthew Hutchinson}
\date{May 2026}

\begin{document}
\maketitle

\section*{Introduction}
This project studies how knowledge editing methods behave beyond the directly edited fact.
The core question is whether an edit remains local and logically consistent when queried through paraphrases, inverse relations, compositional prompts, and short reasoning chains.
I compare three editing paradigms on GPT-2 XL: ROME, MEMIT, and IKE.
ROME and MEMIT modify model weights, while IKE retrieves in-context demonstrations at inference time.

The evaluation center is CounterFact, which provides direct rewrite prompts, paraphrases, and locality prompts.
To extend beyond the standard metrics, I built a 100-probe diagnostic set that targets logical negation, inverse relations, compositional inference, contradiction handling, and short chain-of-thought style prompts.

\section*{Related Work}
ROME \cite{meng2022rome} introduced rank-one model editing by locating and patching a key MLP layer so the model changes one fact while preserving other behavior.
MEMIT \cite{meng2023memit} extends this idea to mass editing by updating multiple facts and layers more efficiently.
IKE \cite{zeng2023ike} takes a different approach: it uses retrieved in-context examples rather than persistent weight updates.

The project follows the EasyEdit implementation of these methods \cite{easyedit}, which provides a common harness for comparing direct edit success, paraphrase transfer, and locality.
That shared harness is useful, but I found that some built-in prompts are noisy, especially for CounterFact rephrase evaluation, so I treat rephrase accuracy as a relative metric across methods rather than a strict paper comparison.

\section*{Methodology}
The experimental setup is GPT-2 XL on CounterFact with EasyEdit.
For ROME and MEMIT, I ran the standard 100-edit single-edit baselines and also the true MEMIT batch sweep at 10, 50, and 100 simultaneous edits.
For IKE, I used the full CounterFact file as the retrieval pool and evaluated the same 100 diagnostic probe cases with retrieved contexts prefixed to the query.

The main metrics are:
\begin{itemize}[leftmargin=*]
  \item \textbf{Rewrite accuracy}: direct success on the edited prompt.
  \item \textbf{Rephrase accuracy}: transfer to a paraphrased prompt.
  \item \textbf{Locality accuracy}: preservation of unrelated prompts.
  \item \textbf{Probe pass rate}: whether a diagnostic prompt produces the expected first token or short generation.
\end{itemize}

The diagnostic probe set is organized into five categories: logical negation, symmetric inverse, compositional, contradiction, and chain-of-thought.
Each probe is also labeled by type:
\texttt{implicit\_edit}, \texttt{target\_conditioned}, or \texttt{supplied\_fact\_reasoning}.
That split matters because supplied-fact prompts are easier than implicit transfer tests.

\section*{Results So Far}
The baseline edit results are stable enough to support a first pass comparison.

\begin{center}
\begin{tabular}{lccc}
\toprule
Method & Rewrite & Rephrase & Locality \\
\midrule
ROME & 1.000 & 0.540 & 0.790 \\
MEMIT & 0.810 & 0.230 & 0.980 \\
IKE & 0.990 & 0.990 & 0.110 \\
\bottomrule
\end{tabular}
\end{center}

These numbers show a clear tradeoff.
ROME is closest to the original paper-level direct rewrite behavior and has decent locality.
MEMIT preserves locality very well in the single-edit setting, but its direct rewrite rate is lower than ROME in this implementation.
IKE achieves strong direct rewrite and paraphrase behavior on the sampled CounterFact subset, but it is much weaker on locality because the in-context examples can perturb unrelated outputs.

The diagnostic probes sharpen the picture.
ROME and MEMIT both improved the total probe pass rate from 36\% pre-edit to 64\% post-edit.
IKE improved from 36\% to 50\%, which is still meaningful but notably lower than the weight-edit methods.
Across all three methods, inverse-style prompts remain the weakest category.
ROME and MEMIT are strongest on logical negation, while IKE is weaker there but does better on some contradiction and supplied-fact reasoning prompts.

\section*{Discussion and Next Steps}
The main lesson so far is that direct rewrite success is not enough to characterize an editing method.
Two methods can both edit the target fact correctly and still behave very differently when the fact is queried indirectly.
The custom probes are useful because they expose exactly that difference.

The next steps are straightforward:
\begin{itemize}[leftmargin=*]
  \item write a concise results section that interprets the probe categories,
  \item run RippleEdits and MQuAKE if time permits,
  \item tighten the narrative around locality failures and inverse-relation failures.
\end{itemize}

\begin{thebibliography}{9}
\bibitem{meng2022rome}
Eric Meng, et al.
\newblock {ROME: Locating and Editing Factual Associations in GPT}.
\newblock \emph{NeurIPS}, 2022.

\bibitem{meng2023memit}
Eric Meng, et al.
\newblock {MEMIT: Mass Editing Memory in a Transformer}.
\newblock \emph{ICLR}, 2023.

\bibitem{zeng2023ike}
Yile Zeng, et al.
\newblock {Knowledge Editing as In-Context Learning}.
\newblock 2023.

\bibitem{easyedit}
EasyEdit project.
\newblock \url{https://github.com/zjunlp/EasyEdit}
\end{thebibliography}

\end{document}
